\section{Conclusion}
\label{sec:conclusion}

% Provisional conclusion (draft 4 Oct 2026): to be rewritten once the hardware reconfiguration
% results are in.
% 9 Oct (structure audit): the findings are now one paragraph per question, in the order of
% Table tab:rq_map (carry, RQ1, RQ2); no finding changed. Still to write: a value statement (what is
% new, who could use it) and research-level future work.

This thesis set out to test whether a centralised controller can carry a cable-suspended payload
through the departure and the mid-flight arrival of individual UAVs, with a centralised load--cable
planner based on \cite{tudelft} and a cable-aware tracking MPC per UAV, flown on Betaflight UAVs over
their radio links with motion capture as the only sensing.

\paragraph{Carrying with a fixed team.} On hardware, three and then four UAVs carried the 0.86\,kg ring
to 1.0\,m, around a circle and through a figure-eight, within 3\,cm of the path on average once the
trackers used the rod pull they measure. The remaining gaps to simulation, a low hover and a steady horizontal offset,
have identified causes or candidate causes. Steady payload errors in the re-planned cascade come
from sources that look alike at the payload (the tracker's relaxed initial condition, a breakaway pull
cap held for the whole flight, model errors and external forces), and each needed a different
remedy. Integral action in the planner's model removed the steady error from external forces in all
three axes in simulation. A per-UAV thrust error is a moment that only a per-UAV correction removes.

\paragraph{RQ1: leave and join.} Reconfiguration can be handled by resizing the planner's optimal
control problem, because the attachment geometry is a runtime parameter. In simulation a detach this
way left the payload level within seconds, and on the rig it held the ring in all four flights,
announced and unannounced, with peak tilts of 6.5--17.8$^\circ$.
\redo{the two comparisons behind this finding: the payload tilt left by a network hand-off on the same detach, and the payload transient of an attach by OCP resize in the full rejoin mission.}
% --- original text, flown before the rig twin; kept to refill this place:
%     the attachment geometry is a runtime parameter. In simulation a detach this way left the payload
%     level within seconds where handing it to a dissipative network left it tilted by
%     26--48$^\circ$. An attach this way, in a full mission with a collaborator's approach controller,
%     rejoined a UAV with a payload transient under 2$^\circ$.
Whether a payload can stay level after a UAV leaves is set by geometry: its centre of mass
must remain inside the polygon of the surviving attachment points, which chooses the layout
(plates 1/3/5/9 for four UAVs).
% --- original text, flown before the rig twin; kept to refill this place:
%     (plates 1/3/5/9 for four UAVs) and explains several earlier failures.

\paragraph{RQ2: trajectory-based against dissipative.}
\redo{the comparison finding: trajectory-based planner against the dissipative network in tracking, detach and attach (Section~\ref{sec:sim_architecture}).}
% --- original text, flown before the rig twin; kept to refill this place:
%     \item On this system the trajectory-based planner outperformed the dissipative network in
%     tracking, detach and attach. The network's advantages did not carry over to rigid rods and a
%     heavier payload.

\paragraph{Future work.} The immediate steps are the attach on the rig and then the full rejoin with
the collaborator's hardware, the model integral through a detach on the rig, and a check of the
motion-capture calibration for the horizontal pull. Beyond this thesis: a model of slack cables,
so the planner can take over before the rods are taut; force sensing at the UAVs, the onboard
equivalent of \cite{tudelft}'s incremental inversion, to remove per-UAV thrust errors; outdoor
flight without motion capture; and larger teams, for which a pre-built solver per team size becomes
a memory and build-time cost.
